% Options for packages loaded elsewhere
\PassOptionsToPackage{unicode}{hyperref}
\PassOptionsToPackage{hyphens}{url}
\documentclass[
]{article}
\usepackage{xcolor}
\usepackage[margin=1in]{geometry}
\usepackage{amsmath,amssymb}
\setcounter{secnumdepth}{-\maxdimen} % remove section numbering
\usepackage{iftex}
\ifPDFTeX
  \usepackage[T1]{fontenc}
  \usepackage[utf8]{inputenc}
  \usepackage{textcomp} % provide euro and other symbols
\else % if luatex or xetex
  \usepackage{unicode-math} % this also loads fontspec
  \defaultfontfeatures{Scale=MatchLowercase}
  \defaultfontfeatures[\rmfamily]{Ligatures=TeX,Scale=1}
\fi
\usepackage{lmodern}
\ifPDFTeX\else
  % xetex/luatex font selection
\fi
% Use upquote if available, for straight quotes in verbatim environments
\IfFileExists{upquote.sty}{\usepackage{upquote}}{}
\IfFileExists{microtype.sty}{% use microtype if available
  \usepackage[]{microtype}
  \UseMicrotypeSet[protrusion]{basicmath} % disable protrusion for tt fonts
}{}
\makeatletter
\@ifundefined{KOMAClassName}{% if non-KOMA class
  \IfFileExists{parskip.sty}{%
    \usepackage{parskip}
  }{% else
    \setlength{\parindent}{0pt}
    \setlength{\parskip}{6pt plus 2pt minus 1pt}}
}{% if KOMA class
  \KOMAoptions{parskip=half}}
\makeatother
\usepackage{graphicx}
\makeatletter
\newsavebox\pandoc@box
\newcommand*\pandocbounded[1]{% scales image to fit in text height/width
  \sbox\pandoc@box{#1}%
  \Gscale@div\@tempa{\textheight}{\dimexpr\ht\pandoc@box+\dp\pandoc@box\relax}%
  \Gscale@div\@tempb{\linewidth}{\wd\pandoc@box}%
  \ifdim\@tempb\p@<\@tempa\p@\let\@tempa\@tempb\fi% select the smaller of both
  \ifdim\@tempa\p@<\p@\scalebox{\@tempa}{\usebox\pandoc@box}%
  \else\usebox{\pandoc@box}%
  \fi%
}
% Set default figure placement to htbp
\def\fps@figure{htbp}
\makeatother
\setlength{\emergencystretch}{3em} % prevent overfull lines
\providecommand{\tightlist}{%
  \setlength{\itemsep}{0pt}\setlength{\parskip}{0pt}}
\usepackage{bookmark}
\IfFileExists{xurl.sty}{\usepackage{xurl}}{} % add URL line breaks if available
\urlstyle{same}
\hypersetup{
  pdftitle={Preprocesamiento de datos},
  pdfauthor={Massielle Coti},
  hidelinks,
  pdfcreator={LaTeX via pandoc}}

\title{Preprocesamiento de datos}
\author{Massielle Coti}
\date{2026-07-16}

\begin{document}
\maketitle

\subsection{Carga de datos}\label{carga-de-datos}

¿Como visualizamos los primeros datos?

When you click the \textbf{Knit} button a document will be generated
that includes both content as well as the output of any embedded R code
chunks within the document. You can embed an R code chunk like this:

\begin{verbatim}
## , , Age = Child, Survived = No
## 
##       Sex
## Class  Male Female
##   1st     0      0
##   2nd     0      0
##   3rd    35     17
##   Crew    0      0
## 
## , , Age = Adult, Survived = No
## 
##       Sex
## Class  Male Female
##   1st   118      4
##   2nd   154     13
##   3rd   387     89
##   Crew  670      3
## 
## , , Age = Child, Survived = Yes
## 
##       Sex
## Class  Male Female
##   1st     5      1
##   2nd    11     13
##   3rd    13     14
##   Crew    0      0
## 
## , , Age = Adult, Survived = Yes
## 
##       Sex
## Class  Male Female
##   1st    57    140
##   2nd    14     80
##   3rd    75     76
##   Crew  192     20
\end{verbatim}

¿Como hacemos para saber de que tipo son los datos (texto/numero)?

\begin{verbatim}
##  'table' num [1:4, 1:2, 1:2, 1:2] 0 0 35 0 0 0 17 0 118 154 ...
##  - attr(*, "dimnames")=List of 4
##   ..$ Class   : chr [1:4] "1st" "2nd" "3rd" "Crew"
##   ..$ Sex     : chr [1:2] "Male" "Female"
##   ..$ Age     : chr [1:2] "Child" "Adult"
##   ..$ Survived: chr [1:2] "No" "Yes"
\end{verbatim}

\subsection{Limpieza de datos}\label{limpieza-de-datos}

¿Cuantos datos faltantes hay para la variable edad (Age)?

\begin{verbatim}
## [1] 177
\end{verbatim}

Si los eliminamos perdemos demasiada informacion ¿Que hacemos?

Solucion: \textbf{Imputarlos} (imputar a la Media o la Mediana)
Recordatorio: La media se ve afectada por los \emph{valores atipicos
(outliers)}

\subsection{Analisis Exploratorio de datos (EDA) con
graficos}\label{analisis-exploratorio-de-datos-eda-con-graficos}

Como determinamos si los ninios (Age \textless{} 12) sobrevivieron mas?
generamos una columna llamada `Child'

\subsection{Nivel 1: R Base (Sin incluir ningun
paquete)}\label{nivel-1-r-base-sin-incluir-ningun-paquete}

Ideal para entender la logica sin dependencias

\subsubsection{Tabla de contigencia
(Frecuencias)}\label{tabla-de-contigencia-frecuencias}

Cuanto hombresy mujeres sobrevivieron?

\begin{verbatim}
##    
##       0   1
##   0 549   0
##   1   0 342
\end{verbatim}

\begin{verbatim}
##    
##       0   1
##   0 100   0
##   1   0 100
\end{verbatim}

\subsection{Box Plot (Edad vs
Supervivencia)}\label{box-plot-edad-vs-supervivencia}

La edad de los que sobrevivieron es similar a la edad de las que no lo
hicieron?
\pandocbounded{\includegraphics[keepaspectratio]{Preprocesamiento-de-datos_files/Preprocesamiento-de-datos_files/figure-latex/unnamed-chunk-6-1.pdf}}

\section{Barplot de supervivencia por
clase}\label{barplot-de-supervivencia-por-clase}

La clase influyo en la supervivencia?
\pandocbounded{\includegraphics[keepaspectratio]{Preprocesamiento-de-datos_files/Preprocesamiento-de-datos_files/figure-latex/unnamed-chunk-7-1.pdf}}
Que porcentaje de cada clase sobrevivio?
\pandocbounded{\includegraphics[keepaspectratio]{Preprocesamiento-de-datos_files/Preprocesamiento-de-datos_files/figure-latex/unnamed-chunk-8-1.pdf}}

\end{document}
